\chapter*{Abstract}
\addcontentsline{toc}{chapter}{Abstract}
\normalsize

Automated facial attendance and access control systems deployed on standard monocular RGB cameras remain vulnerable to presentation attacks, particularly video replays on mobile screens and printed photographs that impersonate enrolled users without specialized depth or infrared sensors. To close this security gap on commodity CPU hardware, this thesis develops a real-time passive face anti-spoofing framework built around a compact $545{,}154$-parameter convolutional neural network (\texttt{LivenessNet}) operating on $32 \times 32$ local facial patches, paired with OpenCV YuNet detection, video-balanced training queues, weight-space interpolation, and a frozen TensorFlow Lite XNNPACK execution runtime. Methodological investigation progressed across two complementary datasets: a multi-source corpus of $28{,}957$ facial crops from $89$ videos used to diagnose normalization drift under imbalanced domain merging, and a dedicated subject-disjoint benchmark of $26{,}094$ facial crops extracted from $91$ 4K recordings on a USB SC-M200 camera. Controlled ablation experiments revealed that replacing Batch Normalization with Group and Layer Normalization alongside stratified $4 \times 8$ source-class-video mini-batches reduced the Attack Presentation Classification Error Rate from $16.16\%$ to $4.44\%$ on the merged benchmark. Subsequent feature-aligned adaptation on the SC-M200 corpus achieved $98.88\%$ validation frame accuracy, $98.09\%$ independent test frame accuracy ($1.94\%$ Average Classification Error Rate, with $99.88\%$ bona fide recall and $99.87\%$ spoof precision), $100.00\%$ video-level classification accuracy across all $14$ unseen test videos, and $99.84\%$ external cross-camera accuracy. Exported to unquantized LiteRT format, the classifier executes in $0.163$~ms median CPU latency, sustaining $30$ frames per second in a multi-threaded camera runtime, $40.61$~ms end-to-end HTTP microservice latency, and reliable spoof blocking within an InsightFace ArcFace web attendance application. These findings confirm that disciplined data partitioning, balanced optimization, and lightweight graph compilation enable practical passive liveness verification on standard office hardware.

\vspace{0.5cm}
\noindent\textbf{Keywords:} Face Anti-Spoofing; Presentation Attack Detection; Lightweight Convolutional Neural Network; Domain Generalization; Biometric Attendance System.
